% ==================================================================
%  THE OR ALGORITHM (pp. 4--5): Pillar 2
% ==================================================================

\section{Decomposition-Guided Design Selection}
\label{sec:algorithm}

The separable structure of the four-term identity transforms agent design
selection from an opaque evaluation problem into a structured optimization.
Once per-mode success probabilities $\{p(s,M_j)\}$ are estimated from a
shared pilot, the four-term objective for \emph{any} design $(M_j, D)$ can
be computed analytically at negligible cost.  This section formalizes the
resulting algorithm and establishes its guarantees.

\subsection{Algorithm}
\label{sec:alg-box}

\begin{figure}[t]
\fbox{\parbox{0.96\textwidth}{\small
\textbf{Algorithm 1: Decomposition-Guided Design Selection}
\label{alg:decomp-guided}

\medskip
\textbf{Input:} Design menu $\cD = \{d_1,\ldots,d_k\}$; task family $\cT$
  with $\Mn$ modes; pilot budget $N_{\mathrm{pilot}}$; confirmation budget
  $N_{\mathrm{conf}}$; candidates~$c$.

\textbf{Output:} Design $d^* \in \cD$ with per-term diagnostics.

\medskip
\textbf{Phase 1 --- Structured Pilot:}
For each mode $s = 1,\ldots,\Mn$ and each model $M_j$ in $\cD$:
run $n_0 = \lceil N_{\mathrm{pilot}}/(m \cdot \Mn)\rceil$ pilot
instances; compute $\hat{p}(s,M_j) = (\text{successes})/n_0$.

\textbf{Phase 2 --- Analytical Scoring} (zero additional evaluations):
For each $d_j = (M_j, D_j, \kappa_j) \in \cD$, compute $\hat{f}(d_j)$
from~\eqref{eq:four-term} using $\hat{p}(\cdot, M_j)$ and known
$\kappa_j, D_j$.  Rank all $k$ designs by $\hat{f}$.

\textbf{Phase 3 --- Confirmation:}
Select top-$c$ candidates $\cC$.  Run each end-to-end on $N_{\mathrm{conf}}$
fresh instances.  Return $d^* = \arg\min_{d \in \cC} \tilde{f}(d)$
with diagnostics $(\hat{\varphi}, \hat{G}, \hat{\mismatch}, \text{cost})$.
}}
\end{figure}

Phase~2 is the key structural insight: it scores all $k$ designs from
shared pilot data using $O(k \cdot \Mn)$ arithmetic operations and zero
task evaluations.  This is what distinguishes Algorithm~1 from benchmark
racing, which must evaluate each design end-to-end.

\subsection{Near-Optimality Guarantee}

\begin{proposition}[Near-optimality]
\label{prop:near-opt}
If the pilot estimates satisfy $|\hat{f}(d) - f(d)| \le \delta$ for every
$d \in \cD$, then the Phase-2 selection $\hat{d}$ satisfies
\begin{equation}
\label{eq:near-opt}
  f(\hat{d}) \;\le\; f(d_{\mathrm{opt}}) + 2\delta,
\end{equation}
where $d_{\mathrm{opt}} = \arg\min_{d \in \cD} f(d)$.
With confirmation error~$\eta$, the final output satisfies
$f(d^*) \le f(d_{\mathrm{opt}}) + 2\delta + 2\eta$.
The bound is tight: a two-design adversarial instance achieves
regret $2\delta - \epsilon$ for any $\epsilon > 0$.
\end{proposition}

\begin{proof}
Since $\hat{d}$ minimizes $\hat{f}$ over $\cD$:
$f(\hat{d}) \le \hat{f}(\hat{d}) + \delta
\le \hat{f}(d_{\mathrm{opt}}) + \delta
\le f(d_{\mathrm{opt}}) + 2\delta$.
The confirmation bound follows by the same argument applied within
$\cC \ni \hat{d}$.  See Appendix~\ref{app:near-opt} for the tightness
construction.
\end{proof}

\subsection{Evaluation Complexity}

\begin{proposition}[Evaluation complexity]
\label{prop:eval-complexity}
Algorithm~1 uses
$N_{\mathrm{Alg1}} = m \cdot \Mn \cdot n_0(\delta) + c \cdot N_{\mathrm{conf}}$
task evaluations, where $m$ is the number of distinct models and
$n_0(\delta) = \log(2\Mn/\dconf) / (2\delta^2 \pmin^2)$.
Benchmark racing uses $N_{\mathrm{race}} = k \cdot n_{\mathrm{eval}}(\delta)$.
The savings factor is
\begin{equation}
\label{eq:savings}
  \frac{N_{\mathrm{race}}}{N_{\mathrm{Alg1}}}
  \;=\; \frac{k \cdot n_{\mathrm{eval}}}{m \cdot \Mn \cdot n_0 + c \cdot N_{\mathrm{conf}}},
\end{equation}
which grows linearly in~$k$ for fixed~$m$ and~$\Mn$.
In particular, $N_{\mathrm{Alg1}}$ is $O(1)$ in~$k$.
\end{proposition}

The constant-in-$k$ property arises because Phase~2 is free: the
separable structure lets all $k$ designs share the same pilot estimates.
Each additional design in the menu adds only $O(\Mn)$ arithmetic, not
task evaluations.

\subsection{Sample Complexity Cascade}

\begin{proposition}[Budget-to-regret cascade]
\label{prop:cascade}
Given a total pilot budget~$N$, with probability at least $1 - \dconf$:
\begin{equation}
\label{eq:cascade}
  f(\hat{d}) - f(d_{\mathrm{opt}})
  \;\le\;
  2\sqrt{\frac{m \cdot \Mn \cdot \log(2\Mn/\dconf)}{2N}}.
\end{equation}
The regret decays as $O(1/\sqrt{N})$ and depends on $m \cdot \Mn$, not
on $k$.  The Bernstein refinement replaces $1/(2N)$ with
$3\pmin/(N)$, improving the bound by $\sqrt{6\pmin}$ when $\pmin \ll 1$.
\end{proposition}

Inverting~\eqref{eq:cascade}: to guarantee regret at most~$R$, the pilot
budget is $N = O(m \cdot \Mn / R^2)$---independent of~$k$.

\subsection{Experimental Validation}
\label{sec:exp10}

We validate on synthetic packed families with $\Mn{=}5$ modes, $m{=}5$ models, and $k \in \{10, 50, 100, 200, 500\}$ designs (50 MC replications each). Table~\ref{tab:exp-combined} shows Algorithm~1 vs four baselines at matched budget (7,500 evaluations).

\begin{table}[t]
\centering
\caption{Algorithm~1 vs baselines at matched budget (7,500 evaluations). Savings factor = racing evals / Alg.\,1 evals. Regret is mean over 50 replications. Algorithm~1's cost is $O(1)$ in $k$; racing and BO degrade as $k$ grows.}
\label{tab:exp-combined}
\small
\begin{tabular}{@{}rrrrrrr@{}}
\toprule
$k$ & Savings & \multicolumn{5}{c}{Mean regret} \\
    & factor  & Alg.\,1 & BO & Racing & SHA & Random \\
\midrule
50  & 3.3$\times$  & 0.045 & 0.049 & 0.057 & \textbf{0.026} & 0.120 \\
100 & 6.7$\times$  & \textbf{0.031} & 0.071 & 0.080 & 0.023 & 0.304 \\
200 & 13$\times$   & \textbf{0.022} & 0.044 & 0.095 & 0.027 & 1.116 \\
500 & \textbf{33$\times$} & \textbf{0.032} & 0.052 & 0.149 & 0.027 & 2.511 \\
\bottomrule
\end{tabular}
\end{table}

The savings factor grows linearly in $k$ as predicted by Proposition~\ref{prop:eval-complexity}, with crossover near $k{\approx}15$. Successive halving is competitive in regret but does not provide per-term diagnostics; at $k{\ge}200$, Algorithm~1 matches its regret while offering interpretable decomposition of why the selected design wins.
